% ---------------------------------------------------
% Glosario de términos usados en Lakdawala et al. (2025)
% Formato: \item[Término] (\textit{término en inglés}) definición
% ---------------------------------------------------

Este glosario reúne los términos técnicos que aparecen en Lakdawala (2025). Entre paréntesis se indica el término en inglés usado en el artículo y, cuando aplica, su notación.

\subsection*{Turbinas y banco de pruebas}

\begin{description}[style=nextline, leftmargin=1.2em, itemsep=3pt, font=\normalfont\bfseries\color{verdeOscuroUN}]

\item[Pico central hidroeléctrica] (\textit{pico hydro}) Instalación hidroeléctrica de muy pequeña potencia, pensada para abastecer cargas aisladas, fuera de la red eléctrica (\textit{off-grid}). El prefijo \textit{pico} se refiere al rango de potencia del generador.

\item[Turbina de impulso] (\textit{impulse turbine}) Turbina en la que el agua sale de una boquilla convertida en un chorro de alta velocidad y transmite su energía cinética al rodete al golpear los cangilones, a presión atmosférica. Las turbinas Pelton y Turgo son de este tipo.

\item[Turbina Pelton] (\textit{Pelton Turbine Runner}, PTR) Turbina de impulso cuyos cangilones tienen dos cavidades separadas por una arista central que divide el chorro en dos mitades, que salen por los costados del cangilón.

\item[Turbina Turgo] (\textit{Turgo Turbine Runner}, TTR) Turbina de impulso con cangilones de una sola cavidad, más angostos que los Pelton, en la que el chorro entra por una cara del rodete y sale por la cara opuesta.

\item[Rodete] (\textit{runner}) Rueda giratoria de la turbina formada por el disco y los cangilones; convierte la energía del chorro en giro del eje.

\item[Cangilón] (\textit{bucket}) Cada una de las cucharas o álabes del rodete donde golpea el chorro. En el estudio se fabricaron por impresión 3D.

\item[Arista divisoria] (\textit{splitter}) Borde central del cangilón que parte el chorro. Su ángulo ($\theta_s$) fue de 30° en el rodete Pelton y de 15° en el Turgo.

\item[Diámetro primitivo] (\textit{pitch circle diameter}, PCD) Diámetro de la circunferencia del rodete sobre la que el chorro golpea los cangilones. Fue de 100~mm en los dos rodetes.

\item[Distancia boquilla-rodete] (\textit{nozzle standoff distance}, $N_{sod}$ o SOD) Distancia entre la punta de la boquilla y el punto donde el chorro golpea el cangilón. Se probaron 40, 50 y 60~mm.

\item[Caída o cabeza] (\textit{head}, $H$) Altura de agua disponible, que determina la energía por unidad de peso del agua. En el estudio se mantuvo constante en 10~m.

\item[Caudal] (\textit{flow rate}, $Q$) Volumen de agua que pasa por la turbina por unidad de tiempo; se midió con un caudalímetro electromagnético.

\item[Velocidad angular] (\textit{angular velocity}, $\omega$) Velocidad de giro del rodete; se midió con un tacómetro sin contacto.

\item[Eficiencia hidráulica] (\textit{hydraulic efficiency}, $\eta_h$) Fracción de la potencia hidráulica del agua que el rodete convierte en potencia en el eje:
\[
\eta_h = \frac{T\,\omega}{\rho\, g\, Q\, H},
\]
donde $\rho$ es la densidad del agua y $g$ la aceleración de la gravedad. Es la variable respuesta de este trabajo.

\item[PLA] (\textit{Polylactic Acid}) Ácido poliláctico, plástico biodegradable usado para imprimir en 3D los cangilones.

\end{description}
